
\section{Experiments}
\subsection{Bind-Your-Avatar-Benchmark}
To the best of our knowledge, no publicly available benchmark has been specifically designed for multi-talking-character scenarios. To comprehensively assess the model's generalizability in animating arbitrary characters in diverse real-world scenarios, we introduce Bind-Your-Avatar-Benchmark, a novel benchmark. We've meticulously collected 40 pairs of portraits from various domains, featuring an eclectic mix of styles, such as Tiefling-like fantasy creatures and sci-fi-inspired humanoid robots. These portraits showcase a wide age range, including adults, young people, teenagers, and infants, set against backgrounds that span from realistic settings to fantastical and science-fiction landscapes.
Each portrait pair is accompanied by corresponding dual-stream speech audio clips, generated by the text-to-speech model \cite{du2024cosyvoice2}, as well as textual prompts. These prompts describe different appearances (e.g., wearing sunglasses, in a black suit), actions (e.g., reach out, head shaking), and environments (e.g., The view out the window is changing). The benchmark comprises two subsets: 20 pairs come with corresponding inpainting frames, while the remaining 20 pairs do not, offering versatility for different research needs and expected to significantly benefit the development of the community.
\subsection{Experimental Setup}
\begin{figure*}[t]
    \centering
    \includegraphics[width=1\textwidth]{figure7} % Reduce the figure size so that it is slightly narrower than the column.
    \caption{
    Qualitative comparison with Ingredients~\cite{fei_ingredients_2025} on the Bind-Your-Avatar-Benchmark (without inpainting frame), which shows that our method have strong ability on both content detail generation and identity preservation.
    }
    \label{fig7}
\end{figure*}
\subsubsection{Implementation Details}
For simplicity, we set the number of characters $n$ to 2 in all experiments. Across all training phases, the three stages consist of 10,000, 40,000, and 10,000 steps, respectively. The total batch size is set to 16, and the learning rate is $1 \times 10^{-5}$. To enhance classifier-free guidance, the reference image, inpainting frame, and audio are each randomly dropped with a probability of 0.05 during training.
The loss weights for $\mathcal{L}_r$, $\mathcal{L}_{st}$, and $\mathcal{L}_{\text{layer}}$ are set to 1, 0.001, and 8, respectively.
During inference, we use 50 sampling steps and set the classifier-free guidance (CFG) scale for both audio and text to 7.
Our training sets include: (1) single-character datasets: Hallo3 \cite{cuiHallo3HighlyDynamic2025}, HDTF \cite{zhang2021flow}, VFHQ \cite{xie2022vfhq}; and (2) multi-character datasets: VICO \cite{zhou2022responsive}, Friends-MMC \cite{wang2025friends}, and our dataset. VFHQ is curated using URL labels due to its silent nature. For VICO, we concatenate per-character videos to form unified multi-character scenes. Friends-MMC is cleaned and preprocessed as with our dataset.

\subsubsection{Evaluation Metrics}
We evaluate all methods based on three aspects: identity preservation, audio-visual synchronization, and visual quality.
\textbf{Identity Preservation:} Following \cite{yuan_identity-preserving_2024,fei_ingredients_2025}, we compute pairwise face similarities (Face Sim) between reference and generated faces using FaceNet \cite{schroff2015facenet}.
\textbf{Audio-Visual Synchronization:} For each character, we use a face detector \cite{yolov8_ultralytics} to crop faces and apply AV\_MossFormer2\_TSE\_16K \cite{clearvoice_2025} for speech separation. We then calculate Sync C and Sync D scores, which are metrics that quantify audio-visual sync quality, by averaging across all characters.
\textbf{Visual Quality:} FID~\cite{heusel2017gans} (on face regions using InceptionV3~\cite{szegedy2016rethinking}) and FVD~\cite{unterthiner2019fvd} are used to evaluate generation quality.
\subsubsection{Test Sets}
For the multi-talking-characters video generation task, we employ two test sets. The first comprises 20 clips derived from our proposed MTCC datasets. The second test set utilizes the 20 pairs from the Bind-Your-Avatar-Benchmark that include inpainting frames.
For the multi-talking-characters video generation (without inpainting frame) task, we use the 20 pairs from the Bind-Your-Avatar-Benchmark that lack inpainting frames. These carefully selected test sets, in conjunction with the benchmark, allow for a thorough and nuanced evaluation of our method across different video generation tasks.
\begin{figure}[t]
    \centering
    \includegraphics[width=0.47\textwidth]{figure6} % Reduce the figure size so that it is slightly narrower than the column. Don't use precise values for figure width.This setup will avoid overfull boxes.
    \caption{Qualitative comparison with other competing methods on the Bind-Your-Avatar-Benchmark (with inpainting frame).}
    \label{fig6}
\end{figure}
\subsection{Comparison with State-of-the-Art}
\subsubsection{Quantitative Evaluation}
To verify the effectiveness of our method, we compare it with several state-of-the-art character animation approaches. As no existing method specifically targets multi-talking-character task, to the best of our knowledge, we select representative baselines from related areas, including traditional talking-head generation methods (Sonic~\cite{jiSonicShiftingFocus2024} and Hallo3~\cite{cuiHallo3HighlyDynamic2025}) and a multi-concept human video generation method (Ingredients~\cite{fei_ingredients_2025}).

Since Hallo3 and Sonic are not designed to handle multi-stream speech audio, we adopt two practical adaptations for implementation: (1) directly using the overlapping speech as input, or (2) providing two separate clean speech streams along with the left and right regions of the inpainting frame as inputs to generate two independent videos, which are then concatenated. The latter approach corresponds to the (Concat) variant illustrated in Figure~\ref{fig5}. Additionally, as Hallo3 and Sonic depend on inpainting frames, they are not included in the multi-talking-character (without inpainting frame) comparison. As for Ingredients, since it does not support audio or inpainting frame inputs, we omit these modalities and append the corresponding driving audio to the generated video for Sync-C and Sync-D evaluation.
\begin{table}[t]
    \centering
    \caption{Quantitative comparisons in multi-talking-characters video generation task on MTCC test set. The best results are in \textbf{Bold}.}
    \label{table2}
    \fontsize{6}{6.5}\selectfont
    \setlength{\tabcolsep}{3pt}
    \renewcommand{\arraystretch}{1.4}
    \begin{tabular}{lC{1.5cm}C{1cm}C{1cm}C{0.9cm}C{0.9cm}}
        \hline\noalign{\hrule height 0.6pt}
        \textbf{Method}
        & \textbf{Face Sim~$\uparrow$(\%)}
        & \textbf{Sync-C~$\uparrow$}
        & \textbf{Sync-D~}
        & \textbf{FID~}
        & \textbf{FVD~} \\
        \hline 
        Ingredients& 48.21&	1.02	&12.22&	122.16	&922.31 \\ 
        Sonic & \textbf{76.92}&	2.25	&9.76	&29.98	&312.12\\ 
        Sonic (Concat) & 75.89&	4.52&	8.89&	36.78&	355.12 \\ 
        Hallo3 & 70.14	&2.15	&10.23&	31.26&	332.68 \\ 
        Hallo3 (Concat)&70.87&	4.12&	8.98	&41.25&	422.63\\ 
        \textbf{Ours} &71.26	& \textbf{4.67}&	\textbf{8.58}&	\textbf{29.55}&	\textbf{308.67} \\ 
        \hline\noalign{\hrule height 1.2pt}
    \end{tabular}
\end{table}
\begin{table}[t]
    \centering
    \caption{Quantitative comparisons in multi-talking-characters video generation on Bind-Your-Avatar-Benchmark. The best results are in \textbf{Bold} and the second best results are \underline{underlined}.}
    \label{table3}
    \fontsize{7.5}{6.5}\selectfont
    \setlength{\tabcolsep}{3pt}
    \renewcommand{\arraystretch}{1.4}
    \begin{tabular}{lC{1.77cm}C{1.77cm}C{1.77cm}}
        \hline\noalign{\hrule height 0.6pt}
        \textbf{Method}
        & \textbf{Face Sim~$\uparrow$(\%)}
        & \textbf{Sync-C~$\uparrow$}
        & \textbf{Sync-D~} \\
        \hline
        Ingredients & 51.22	&0.95	&12.97 \\ 
        Sonic &\textbf{71.50}	&2.32	&9.82  \\ 
        Sonic (Concat) & 64.78	&\textbf{4.39}&	8.89  \\ 
        Hallo3 & 64.67	&2.12	&10.55  \\ 
        Hallo3 (Concat) & 61.29	&3.79&	8.72\\ 
        \textbf{Ours} & \underline{65.78}	&\underline{4.09}	&\textbf{8.68} \\ 
        \hline\noalign{\hrule height 0.6pt}
    \end{tabular}
\end{table}
% \begin{figure*}[t]
%     \centering
%     \includegraphics[width=1.01\textwidth]{figure1} % Reduce the figure size so that it is slightly narrower than the column.
%     \caption{
%     Ov
%     }
%     \label{fig5}
% \end{figure*}
Quantitative comparisons on both the multi-talking-character task and the multi-talking-character (without inpainting frame) task are presented in Table~\ref{table2},~\ref{table3}, and Table~\ref{table4}, respectively. Our method outperforms most baselines across the majority of metrics, demonstrating superior audio-visual synchronization and identity preservation.
\subsubsection{Qualitative Evaluation}
To demonstrate the visual effectiveness of our method, we compare and visualize the results alongside several competitive approaches, as shown in Figure~\ref{fig5} and Figure~\ref{fig6} for the multi-talking-character task, and in Figure~\ref{fig7} for the multi-talking-character (without inpainting frame) task. The red highlighted box in Figure~\ref{fig6} indicates a failure case of the competing methods in responding to two-stream speech input and generating a unified background, which our method successfully addresses. Moreover, our method is capable of producing dynamic and vivid backgrounds—with natural lighting variations and scene dynamics such as moving pedestrians in the distance and subtle light changes on the speaker’s face—whereas other state-of-the-art methods tend to produce static and less expressive scenes. This advantage stems from the diverse, in-the-wild multi-speaker scenarios provided by our proposed MTCC dataset.
As shown in Figure~\ref{fig7}, our method not only effectively captures environment details and spatial relationships between characters as described in the prompt, but also maintains strong identity preservation and precise correspondence between each character and their respective speech stream.
% \subsubsection{Multi-Concept Audio-Based Human Video Generation}
\begin{table}[t]
    \centering
    \caption{Quantitative comparisons in multi-talking-character (without inpainting frame) task on Bind-Your-Avatar-Benchmark.}
    \label{table4}
    \fontsize{7.5}{6.5}\selectfont
    \setlength{\tabcolsep}{3pt}
    \renewcommand{\arraystretch}{1.4}
    \begin{tabular}{lC{1.91cm}C{1.91cm}C{1.91cm}}
        \hline 
        \textbf{Method}
        & \textbf{Face Sim~$\uparrow$(\%)}
        & \textbf{Sync-C~$\uparrow$}
        & \textbf{Sync-D~} \\
        \hline
        Ingredients & 48.2 & 1.02 & 12.22  \\ 
        \textbf{Ours} & \textbf{68.9} & \textbf{4.52} & \textbf{7.89}  \\ 
        \hline\noalign{\hrule height 0.6pt}
    \end{tabular}
\end{table}

% \subsection{Evaluations and Comparisons}
% \subsubsection{Evaluation:}The generated video is evaluted based on the following criteria:
% \begin{enumerate}
%     \item \textbf{Identity Preservation:} Each generated character should accurately preserve the identity and facial features of the corresponding reference images.
%     \item \textbf{Audio-Visual Synchronization:} Mouth movements, facial expressions, and head poses of each character should be precisely synchronized with their respective speech audio.
%     \item \textbf{Multi-Character Interaction:} Characters should interact naturally within the same scene, exhibiting contextually appropriate actions, expressions, and gaze, reflecting the dynamics of real conversations.
%     \item \textbf{Visual Quality:} The generated video should exhibit high visual fidelity, including realistic facial animations, smooth transitions, and minimal artifacts.
%     \item \textbf{Text-Condition Alignment:} The actions, expressions, and scene composition should be consistent with the provided text prompt, ensuring narrative coherence and relevance.
% \end{enumerate}
% \end{subsection}
\subsection{Ablation Studies}
% w/w.o router(后面补量化，这次不做了）
\subsubsection{Bounding Box vs Fine-grained Mask}
To investigate whether our proposed fine-grained mask yields more precise results, we compare its performance with traditional bounding box masks in close-interaction scenes, a common practice in talking-head methods~\cite{cuiHallo2LongDurationHighResolution2024}. Figure~\ref{fig8}(a) demonstrates cases where our fine-grained mask outperforms bounding boxes. In close interactions, overlapping bounding boxes fail to correctly bind audio to characters, while our dense masks accurately separate facial regions for precise correspondence.
% w/w.o 3d mask 开车场景
% w/w.o refined mask 老人和小孩
% only qualitative
\begin{figure}[t]
    \centering
    \includegraphics[width=0.49\textwidth]{figure8} % Reduce the figure size so that it is slightly narrower than the column. Don't use precise values for figure width.This setup will avoid overfull boxes.
    \caption{Ablation on (a) bounding box vs fine-grained mask, and (b) 2D mask vs 3D mask.}
    \label{fig8}
\end{figure}
\subsubsection{2D Mask vs 3D Mask}  
To evaluate the role of spatiotemporal mask adaptability in dynamic scenes, we contrast static 2D masks (extracted from the inpainting frame and fixed across frames) with our dynamic 3D masks. As shown in Figure \ref{fig8}(b), in a driving scenario with character motion, the 2D mask—anchored to the initial frame—fails to track the driver’s posture changes. This misaligns the mask with the actual face in subsequent frames, causing audio-character binding errors.
In contrast, our 3D mask dynamically updates frame-by-frame, aligning with the driver’s changing position to maintain precise audio-visual synchronization. This highlights that 3D masks are critical for dynamic scene modeling.